\documentclass[10pt,a4paper]{amsart}
\usepackage[margin=19mm]{geometry}
\usepackage{amsmath,amssymb,mathtools}
\usepackage[scr=rsfs,cal=euler]{mathalfa}
\usepackage{xfrac}
\usepackage[unicode,hidelinks,bookmarks=false]{hyperref}
\newcommand{\ie}{i.e.\ }
\newcommand{\cf}{cf.\ }
\newcommand{\resp}{respectively\ }
\newcommand{\Subsection}{\subsection}
\providecommand{\nobreakdash}{\nobreak\hbox{-}}
\setlength{\emergencystretch}{2em}
\multlinegap0pt
\usepackage{amssymb}
\usepackage{dsfont}
\newtheorem{lemma}{Lemma}
\newtheorem{proposition}{Proposition}
\newcommand{\Loc}{\operatorname{Loc}}
\newcommand{\Spc}{\operatorname{Spc}}
\newcommand{\cat}{\mathcal}
\newcommand{\id}{\mathrm{id}}
\newcommand{\iso}{\xrightarrow{\raisebox{-.6ex}[0ex][0ex]{$\scriptstyle{\sim}$}}}
\newcommand{\loc}{\operatorname{loc}}
\newcommand{\supp}{\operatorname{supp}}
\title{Tensor triangular geometry\\ Notes for an Oberwolfach Seminar}
\author{Henning Krause}
\date{Source: arXiv:2605.19983v1 (CC BY 4.0)}
\begin{document}
\maketitle

\noindent\emph{Source and licence.} Tensor triangular geometry\\ Notes for an Oberwolfach Seminar. Version arXiv:2605.19983v1 (CC BY 4.0), distributed by arXiv under the Creative Commons Attribution 4.0 International licence, http://creativecommons.org/licenses/by/4.0/, as recorded in the arXiv licence metadata for this exact version and shown on the abstract page https://arxiv.org/abs/2605.19983v1. Full licence notice: \texttt{licenses/arxiv-2605.19983-CC-BY-4.0.txt}.\par\smallskip

\noindent\textit{Editorial scope.} Counted unit: the article's proposition pr:stratification-tensor (source0.tex lines 2731--2741). The article's complete original proof of it is delivered with it (source0.tex lines 2743--2782, one \texttt{proof} environment of 40 lines). No mathematics is written, repaired or re-derived: the statement and the proof are the article's own lines, in the article's own notation, and no retained line is substituted or re-flowed.  Retained uncounted as the source-local context the proof needs: lines 2674--2677.  Nothing was omitted from any retained span, so the package carries no omission record.  No retained line contains a citation, so the wrapper carries no bibliography and no bibliography entry.  No retained line contains a figure environment, an image-inclusion command or drawing code, so no image is delivered and the package declares no asset dependency.  Editorial formatting only, in addition: an AMS \texttt{amsart} preamble that restates the article's own declarations and macros used by the retained lines, namely the environments \texttt{lemma}, \texttt{proposition} on separate counters, together with the standard packages those lines need.  The retained environments print in this document as Lemma 1, Proposition 1 in order of appearance, whereas the article prints the same items with the numbering of its own full document.  The complete native source of the article as distributed in the arXiv e-print for this exact version is bundled unchanged as \texttt{sources/200957/source0.tex}.
\begin{lemma}\label{le:tensor-supp-surjective}
  For each subset $U\subseteq \Spc\cat T^{\mathrm c}$ we have
  $\supp_\otimes(\tau_\otimes(U))=U$.
\end{lemma}

\begin{proposition}\label{pr:stratification-tensor}
For a rigidly-compactly generated tensor triangulated category $\cat T$, the
following are equivalent.
\begin{enumerate}
\item The category $\cat T$ is stratified via $\Spc\cat T^{\mathrm c}$.
\item The map $\supp_\otimes$ induces
  a lattice isomorphism $\Loc_\otimes(\cat T)\iso P(\Spc\cat T^{\mathrm c})$.
\item $\supp_\otimes(x)\subseteq\supp_\otimes(y)$ implies
  $\loc_\otimes(x)\subseteq\loc_\otimes(y)$ for all $x,y\in\cat T$.
\end{enumerate}
\end{proposition}

\begin{proof}
  (1) $\Rightarrow$ (2) The assignment
  $ U\mapsto \bigvee_{p\in U}\cat T_{\{p\}}$ yields an inclusion
  preserving map $P(\Spc\cat T^{\mathrm c})\to \Loc_\otimes(\cat T)$ which
  satisfies $\supp_\otimes(\bigvee_{p\in U}\cat T_{\{p\}})=U$.
  Assuming that $\cat T$ is stratified, this map is an inverse for
  $\supp_\otimes$, since for any localising tensor ideal
  $\cat S\subseteq\cat T$ one has $\Gamma_p\cat S=\cat T_{\{p\}}$ for
  each $p\in\supp_\otimes(\cat S)$. Thus $\supp_\otimes$ provides a
  lattice isomorphism $\Loc_\otimes(\cat T)\iso P(\Spc\cat T^{\mathrm c})$.

  (2) $\Rightarrow$ (1) Suppose that $\supp_\otimes$ is a
  bijection. For a localising tensor ideal $\cat S$ we have
\[\supp_\otimes(\cat S)=\supp_\otimes\left(\bigvee_{p\in\Spc\cat T^{\mathrm c}}\Gamma_p\cat S\right)\]
and therefore
\[\cat S=\bigvee_{p\in\Spc\cat T^{\mathrm c}}\Gamma_p\cat S\,.\]
Thus the local-global principle holds for localising tensor
ideals. For a nonzero localising tensor ideal
$\cat S\subseteq \cat T_{\{p\}}$ we have
$\supp_\otimes(\cat S)=\{p\}$, and therefore $\cat S=\cat T_{\{p\}}$. We
conclude that $\cat T$ is stratified.

(2) $\Rightarrow$ (3) This is clear since
$\supp_\otimes(\loc_\otimes (x))=\supp_\otimes(x)$ for all
$x\in\cat T$.

  (3) $\Rightarrow$ (2) We need to show that
  $\supp_\otimes\colon\Loc_\otimes(\cat T)\to P(\Spc\cat T^{\mathrm c})$ is
  injective. Then this map is bijective since
  $\supp_\otimes\tau_\otimes=\id$ by
  Lemma~\ref{le:tensor-supp-surjective}. Suppose that $\cat S,\cat S'$ is a
  pair of localising tensor ideals such that
  $\supp_\otimes(\cat S)= \supp_\otimes(\cat S')$, but
  $\cat S\neq\cat S'$. Choose $x\in\cat S\smallsetminus\cat S'$. Then
  $\supp_\otimes(x)\subseteq
  \bigcup_i\supp_\otimes(x_i)=\supp_\otimes(\coprod_i x_i)$ for some
  family of objects $x_i\in\cat S'$. This implies
  $\loc_\otimes(x)\subseteq\loc_\otimes(\coprod_i x_i)\subseteq\cat
  S'$, which is a contradiction. Thus $\supp_\otimes$ is injective.
\end{proof}

\end{document}
